% Generated draft: requires mathematical and source review.
\documentclass[11pt]{article}
\usepackage[margin=25mm]{geometry}
\usepackage{fontspec}
\setmainfont{Latin Modern Roman}
\usepackage{amsmath,amssymb,mathtools,unicode-math}
\setmathfont{Latin Modern Math}
\usepackage{enumitem,needspace,xurl,hyperref}
\setlength{\parindent}{0pt}
\setlength{\parskip}{5pt}
\setlength{\emergencystretch}{3em}
\newcommand{\N}{\mathbb{N}}
\newcommand{\Z}{\mathbb{Z}}
\newcommand{\Q}{\mathbb{Q}}
\newcommand{\R}{\mathbb{R}}
\newcommand{\C}{\mathbb{C}}
\newcommand{\F}{\mathbb{F}}
\newcommand{\A}{\mathbb{A}}
\newcommand{\PP}{\mathbb{P}}
\newcommand{\E}{\mathbb{E}}
\newcommand{\eps}{\varepsilon}
\newcommand{\dd}{\,\mathrm{d}}
\begin{document}
\section*{UnsolvedMath ID: 800594}
\textbf{Source problem code:} catalogue entry
\subsection{Definitions and mathematical statement}
Let \(X\) be a topological space. A cover of \(X\) is a family \(\mathcal U\) of subsets of \(X\) such that \(\bigcup\mathcal U=X\) and \(X\notin\mathcal U\). It is open if every member is open in \(X\). Following the supplied definition, \(\mathcal U\) is a large cover if \(\{U\in\mathcal U:x\in U\}\) is infinite for every \(x\in X\). It is a \(\tau\)-cover if it is large and, for every \(x,y\in X\), at least one of the sets
\[
\{U\in\mathcal U:x\in U,\ y\notin U\},\qquad
\{U\in\mathcal U:y\in U,\ x\notin U\}
\]
is finite.

Using the standard interpretations of the symbols not defined in the supplied excerpt, an \(\omega\)-cover is a cover \(\mathcal U\) such that every finite subset \(F\subseteq X\) is contained in some \(U\in\mathcal U\). A \(\gamma\)-cover is an infinite cover \(\mathcal U\) such that \(\{U\in\mathcal U:x\notin U\}\) is finite for every \(x\in X\). Write \(\Omega(X)\), \(\Gamma(X)\), and \(\mathrm T(X)\) for the respective collections of admissible open \(\omega\)-, \(\gamma\)-, and \(\tau\)-covers. The supplied excerpt does not specify whether admissible covers are all open covers or only countable open covers; this convention remains to be recovered from the source context. Under a consistent convention, the source states
\[
\Gamma(X)\subseteq\mathrm T(X)\subseteq\Omega(X).
\]

For collections \(\mathcal A\) and \(\mathcal B\) of covers of \(X\), interpret \(\binom{\mathcal A}{\mathcal B}\) as the property
\[
\forall\mathcal U\in\mathcal A\ \exists\mathcal V\subseteq\mathcal U
\quad \mathcal V\in\mathcal B.
\]
Thus the selected family is a subfamily that still covers \(X\).

The intended question, under these standard interpretations, is whether the two properties \(\binom{\Omega}{\Gamma}\) and \(\binom{\Omega}{\mathrm T}\) hold for exactly the same spaces in the source's intended class \(\mathscr C\). Explicitly, is it true that, for every \(X\in\mathscr C\),
\[
\begin{aligned}
&\left[\forall\mathcal U\in\Omega(X)\ \exists\mathcal V\subseteq\mathcal U:
\mathcal V\in\Gamma(X)\right]\\
&\qquad\Longleftrightarrow
\left[\forall\mathcal U\in\Omega(X)\ \exists\mathcal V\subseteq\mathcal U:
\mathcal V\in\mathrm T(X)\right]?
\end{aligned}
\]
The supplied inclusion identifies the unresolved direction as whether \(\binom{\Omega}{\mathrm T}\) implies \(\binom{\Omega}{\Gamma}\). The class \(\mathscr C\), including any separation, metrizability, or restriction to sets of reals, is not specified in the supplied record. Consequently, this is a conditional reconstruction of the formulation, pending recovery of that domain and the cover convention; it does not assert that arbitrary topological spaces are the intended domain.
\subsection{Short English statement}
Does the property that every open \(\omega\)-cover contains an open \(\tau\)-subcover imply that every open \(\omega\)-cover contains an open \(\gamma\)-subcover? The intended class of spaces and any countability restriction on covers are missing from the supplied excerpt. The registry labels the problem open; this is supplied source status, not an independently verified literature assessment.
\subsection{Sources}
\begin{itemize}
\item Local UnsolvedMath registry, permanent ID \(800594\), exact source problem code \texttt{catalogue entry}; title: \emph{Selection Principles in Mathematics --- Problem 21 ([spm])}. The record credits Tsaban and gives proposed year 2006. Its supplied status is \texttt{open}.
\item Supplied source list: Tsaban, \emph{Selection Principles and Special Sets of Reals: Open Problems} (2006), Problem 21. Supplied source URL: \url{https://arxiv.org/abs/math/0609601}. The record identifies the extraction method as \texttt{arxiv\_latex} and reports an access date of 2026-10-16. This is supplied provenance, not a claim of independent source access or verification.
\item The supplied background gives attribution/context as [spm], and the definition of a \(\tau\)-cover carries reference label [tau]. Bibliographic entries resolving these labels were not supplied.
\end{itemize}
\end{document}
